\documentclass[10pt]{article}
\usepackage[margin=20mm]{geometry}
\usepackage{amsmath,amssymb,mathtools,fontspec,unicode-math,xurl,hyperref}
\IfFontExistsTF{DejaVu Serif}{\setmainfont{DejaVu Serif}}{\setmainfont{Latin Modern Roman}}
\setmathfont{latinmodern-math.otf}
\setlength{\parindent}{0pt}\setlength{\parskip}{6pt}\sloppy
\title{Initial scientific review of
Distortion of Hausdorff measures under
Orlicz–Sobolev maps}
\author{}
\date{Frozen batch M25-0247; initial review, decision pending}
\begin{document}
\maketitle
Abstract: The assessment is Major. Equation (4.20) supplies a false identity to the gauge monotonicity proof used in the main image-cover estimate. Separately, the Cantor construction in Theorem 6.2 prescribes incompatible cardinalities and cylinder masses, while the energy estimate requires an admissible construction with that mass control. These invalid derivations affect the distortion and sharpness contributions. The report supplies corrected gauge and Cantor arguments that preserve the claimed conclusions within the checked scope. Five other findings concern the one-dimensional constant, exponential inverse comparison, limiting direction and normalization. No counterexample to Theorems 3.1, 3.4, 3.5, 3.7, 6.1 or 6.2 is established. Repairability does not lower the severity of the printed core derivations. Scientific decision and adjudication remain pending.\par
Article. Andrea Cianchi, Mikhail V. Korobkov and Jan Kristensen. Journal identifier: 10.4171/JEMS/1620.\par
Version examined. The 37-page manuscript `arXiv:2208.08152v2`, 26 May 2023. PDF and printed page numbers coincide. The publisher records online publication on 19 March 2025. The frozen 2025 record is retained; equivalence between this manuscript and the journal version has not been established.\par
\section*{Scope and central mechanism}
The primary reviewer read all 37 source-text pages and actually inspected all 35 rendered body/support pages. Page 28 was recovered through a complete UTF-8 read with ASCII-escaped JSON after the text reader rejected it; the actual PDF page was also checked. They did not perform source-image inspection. Roles and model self-reports are recorded separately from the controller’s UI observation; this report is an initial review, not an independent acceptance.\par
The main construction uses \[J(s)=sB^{-1}\!\left(\frac{\varphi(s)}{s^n}\right),
 \qquad \psi(r)=\varphi(J^{-1}(r)).\] The Young function \(B\) is obtained from \(A\) by (2.27), with \(B=A\) when \(n=1\). The restrictions (2.26), (2.29), (2.30) give finite positive functions, suitable endpoint behavior, and classical inverses. The standing condition \(\varphi(s)/s^n\) nonincreasing must remain part of every gauge statement. The crossing inequality in Lemma 4.4 combines with the Orlicz–Morrey estimate (2.28) and disjoint dyadic covers. The stronger finite-measure-to-zero conclusion additionally uses \(\varphi(s)/s^n\to\infty\) and the two small-scale conditions (3.10).\par
The proof was checked at the algebraic, covering, modular-integral and limiting stages. The random-series sharpness argument was checked through the cutoff norm, branching, overlap, coefficient summability and energy calculation. The external Orlicz–Morrey theorem with its precise constant, and the full general theorem on continuity from below of Hausdorff content, were not independently reproved. No additional supplement or code artifact was acquired in the bounded search. The publisher’s PDF required a subscription, and an Oxford repository request returned HTTP 403; neither limitation was bypassed or classified as a mathematical defect.\par
\section*{Confirmed issues}
\subsection*{F1. The introductory positive-measure assertion is too strong (Minor)}
Location. Page 3, the final sentence of the first paragraph.\par
For every \(0<\alpha\le n\), \[\mathcal H^\alpha(E)=0\quad\Longrightarrow\quad \mathcal L^n(E)=0.\] For \(\alpha<n\), apply \(\operatorname{diam}(U)^n\le\delta^{n-\alpha}\operatorname{diam}(U)^\alpha\) to sufficiently fine covers of bounded pieces; the endpoint \(\alpha=n\) follows from proportionality with Lebesgue measure. The Lusin \(N\) property stated in the same paragraph therefore forces \(\mathcal H^n(u(E))=0\) for the indicated target \(\mathbb R^n\).\par
Already \(n=m=1\) supplies a direct obstruction independent of the higher-dimensional cited result: the relevant space is locally absolutely continuous, and such maps send length-zero sets to length-zero sets. Thus the printed assertion about a positive \(n\)-dimensional image measure cannot hold. An image can have Hausdorff dimension \(n\) while having zero \(\mathcal H^n\) measure. Retain the supported full-dimension assertion and remove the positive-measure inference. This concerns the introduction, not a counterexample to the main distortion results.\par
\subsection*{F2. Exact normalization must be separated from multiplicative equivalence (Minor)}
Location. Page 5, the assertion preceding (2.6); p. 10, Remark 3.2; related power normalization in (3.24) on p. 13.\par
In dimension one, \(\varphi(r)=2r\) gives \[\mathcal H^\varphi=2\mathcal L^1\ne\mathcal L^1,\qquad
 \frac{\varphi(r)}r=2.\] Thus literal inequality of the measures is not equivalent to the divergent limit in (2.6). The statement should concern failure of equivalence up to multiplicative constants. The actual divergent-limit hypothesis used in Theorem 3.4 is unaffected.\par
Likewise, if \(\varphi(s)=s^n\), the definitions give exactly \[J(s)=sB^{-1}(1),\qquad
 \psi(r)=\frac{r^n}{[B^{-1}(1)]^n}.\] For \(n=1\) and \(A(t)=4t^2\), this is \(\psi(r)=2r\), not \(r\). One may retain the constant, use multiplicative equivalence, or explicitly normalize \(B(1)=1\). For \(n\ge2\) and \(A=t^p\), direct conjugation gives \[B(t)=c_Bt^p,\quad c_B=(n'-p')^{p-1},\qquad
 \psi(r)=c_B^{\alpha/(p+\alpha-n)}
 r^{\alpha p/(p+\alpha-n)}.\] Thus the exact equality in (3.24) also needs its normalization. Positive scalar factors preserve zero sets, finiteness and the distortion exponent, which explains the limited impact.\par
\subsection*{F3. The explicit constant in Example 3.6 fails when \(n=1\) (Minor)}
Location. Page 12, equation (3.16).\par
Take \[n=m=1,\quad p=2,\quad\alpha=1,\qquad
 u(x)=\min\{\max\{x,0\},1\},\quad E=[0,1].\] The homogeneous Sobolev assumption is satisfied, and \[\|u'\|_{L^2(\mathbb R)}=1,\qquad
 \mathcal H_\infty^1(E)=\mathcal H_\infty^1(u(E))=1.\] Under the usual convention \(n'=\infty\), the factor \[\left(\frac{p-1}{n'-p'}\right)^{
       \alpha(p-1)/(\alpha+p-n)}\] vanishes. Equation (3.16) would therefore give \(1\le0\). Without that convention, the formula is undefined for an admitted dimension.\par
Complete one-dimensional replacement. Restrict the printed \(n'-p'\) formula to \(n\ge2\). When \(n=1\), put \[\tau=\frac{\alpha p}{p+\alpha-1},\qquad 0<\alpha\le1,\quad p>1.\] On an interval \(I\), \[\operatorname{diam}u(I)\le |I|^{1-1/p}
 \left(\int_I|u'|^p\right)^{1/p}.\] Replace covering sets by their convex intervals and merge overlapping intervals. Since \(\alpha\le1\), this does not increase the sum of their \(\alpha\)-powers. The remaining intervals have disjoint interiors. Hölder’s inequality for the resulting sum gives \[\sum_I\operatorname{diam}u(I)^\tau
 \le \|u'\|_p^\tau
       \left(\sum_I |I|^\alpha\right)^{(p-1)/(p+\alpha-1)}.\] Infimizing proves \[\mathcal H_\infty^\tau(u(E))
 \le \|u'\|_p^\tau
       [\mathcal H_\infty^\alpha(E)]^{(p-1)/(p+\alpha-1)}.\] The clamp example attains equality for \(\alpha=1\). Theorem 3.5 itself uses \(B=A\) in dimension one and is not refuted by the bad specialized constant.\par
\subsection*{F4. Equation (4.20) omits the \(n\)-th power (Major)}
Location. Page 19, Lemma 4.3 and its following monotonicity argument.\par
With \(v=\varphi(s)/s^n\), the correct identity is \[\frac{\varphi(s)}{J(s)^n}
 =\frac{v}{[B^{-1}(v)]^n}.\] The source prints the denominator without its \(n\)-th power. For an exact test, take \[n=2,\quad A(t)=\frac{27}{8}t^4,\quad B(t)=t^4,\quad
 \varphi(s)=s,\quad s=16.\] The value of \(B\) follows directly from (2.27), or from the conjugation formula in F2. Then \(J(16)=8\), so the left side is \(1/4\), whereas the printed right side is \((1/16)/(1/2)=1/8\).\par
The proof also needs the corresponding stronger monotonicity. Namely, \[\frac{v}{[B^{-1}(v)]^n}=\frac{B(z)}{z^n},
 \qquad z=B^{-1}(v),\] is nondecreasing by the already established (4.6). Since \(v(s)\) is nonincreasing, the composition is nonincreasing, exactly proving (4.17). In dimension one the same assertion follows from convexity of \(A=B\). This corrected argument establishes the gauge lemma; its conclusion is not contradicted.\par
Dependency and scope. The source proves (4.6), but the actual Lemma 4.3 proof uses the false (4.20) and the different monotonicity of \(v/B^{-1}(v)\). The available premise supplies a replacement justification, not a second printed proof of (4.17). On p. 23, Theorem 3.1 invokes Lemma 4.3 in (5.6) to convert normalized image-cover costs to unscaled costs; this gives (5.9) and the measure conclusion and is retained in (5.13) for the content conclusion. Lemma 4.10 on pp. 22–23 also invokes (4.17) for the monotonicity of \(\Xi_{\psi,B}\). The invalid derivation therefore has a demonstrated core dependency, warranting Major classification. The corrected calculation preserves those uses; no false main theorem is claimed.\par
\subsection*{F5. The critical exponential inverse comparison is false (Minor)}
Location. Page 14, the \(\alpha=0\) branch of (3.28). Here and below, \(\asymp\) denotes comparison by positive multiplicative constants.\par
Take \(n=1\), \(q=1\), \(\beta=-1\), and \[A(t)=t\log(1+t).\] This is a finite positive Young function with \[A''(t)=\frac{t+2}{(t+1)^2}>0,\qquad
 \lim_{t\to0}\frac{A(t)}t=0,\quad
 \lim_{t\to\infty}\frac{A(t)}t=\infty.\] For \(L=\log(1/s)\ge16\), define \[\varphi(s)=sA\!\left(\frac4{sL^2}\right)
 =\frac4{L^2}\log\!\left(1+\frac{4e^L}{L^2}\right).\] Extend it beyond \(s=e^{-16}\) linearly, with \(\varphi(s)/s\) constant. This is an admissible increasing continuous gauge with \(\varphi(s)/s\) nonincreasing and \(\varphi(s)\asymp1/\log(1/s)\). For completeness, if \(z=4e^L/L^2\), its derivative with respect to \(L\) has sign \[-2\log(1+z)+(L-2)\frac{z}{1+z}
 <-L+4\log L-2\log4-2<0
 \quad(L\ge16).\] The ratio \(\varphi(s)/s=A(4e^L/L^2)\) decreases with \(s\), including across the linear extension.\par
Since \(B=A\), the definitions give exactly \[J(s)=\frac4{[\log(1/s)]^2},\qquad
 J^{-1}(r)=e^{-2/\sqrt r}.\] The printed comparison in (3.28) is \(e^{-1/\sqrt r}\). Their ratio is \(e^{-1/\sqrt r}\to0\), so they are not comparable in the stated sense.\par
Counterevidence and repair. The final gauge nevertheless satisfies \[\psi(r)=r\log(1+re^{2/\sqrt r}),\qquad
 \frac{\psi(r)}{2\sqrt r}\longrightarrow1.\] Thus the final exponent in (3.26) survives. In general, replace the first branch of (3.28) by \[\log\frac1{J^{-1}(r)}
 \asymp r^{-n/(q-n+1-\beta)},\] or by two exponential bounds retaining constants inside their exponents. Composing with the logarithmic \(\varphi\) then gives the claimed final power. The displayed \(J(s)\) asymptotics on pp. 14–15 concern \(s\to0\), so their near-infinity labels should also be corrected. This is a false intermediate comparison with an explicit corrected derivation, not a refutation of Example 3.9’s final conclusion.\par
\subsection*{F6. The zero-volume limit and image labels in Theorem 3.7 need correction (Minor)}
Location. Pages 27–28, especially (5.28) and the following paragraph.\par
When \(\mathcal L^n(E)=0\), the bound involves \[\frac{[B^{-1}(r)]^n}{r}=\Phi_B(1/r).\] It tends to zero as \(r\to\infty\), not as \(r\to0\). For \(n=1\), \(A(t)=4t^2\), it equals \(1/(2\sqrt r)\) and diverges at zero. Reverse the printed limiting direction.\par
The cube cover directly controls \[\mathcal H_\infty^n(u(E))\le\sum_i\operatorname{diam}(u(Q_i))^n\] in \(\mathbb R^m\). The printed \(\mathcal L^n(u(E))\) on pp. 27–28 should retain this content label, especially since \(m>n\) is allowed. With \(V=\mathcal L^n(E)>0\), the choice \[r=\frac{\kappa}{n^{n/2}V}\] in the corrected cover estimate gives exactly \[2\kappa^{n+1}\|\nabla u\|_{L^A}^n
 \Phi_B\!\left(\frac{n^{n/2}V}{\kappa}\right).\] For \(V=0\), use \(r\to\infty\). Exhaustion and content continuity from below give the general-set conclusion. These corrections complete the intended step. The zero-content case of Theorem 3.5 can similarly cite Theorem 3.1(ii), which also covers \(\varphi(r)=r^n\), instead of relying on the more restricted Theorem 3.4.\par
\subsection*{F7. The Cantor construction data are inconsistent as printed (Major)}
Location. Page 32, (6.12)–(6.15), and their use on pp. 33–34.\par
The source specifies \(\#\mathcal A_j=\lfloor2^{j\nu}\rfloor\) and \(0<\mathcal H^\varphi(M)<\infty\), but (6.14) gives cylinder mass at most \(c2^{-j}\). The allowed value \(\nu=1/2\) yields \[0<\mathcal H^\varphi(M)
 \le c\lfloor2^{j/2}\rfloor2^{-j}
 \le c2^{-j/2}\longrightarrow0,\] a contradiction. The required and subsequently used bound on p. 34 is \(c2^{-j\nu}\).\par
There is a separate finite-level obstruction. For \(n=1,\nu=3,j=1\), the construction asks for eight different dyadic intervals of length \(1/4\) in \([0,1]\), where only four exist. The theorem permits these parameters.\par
A construction closing both defects. Start at sufficiently large \(j_0\) and set \[\delta_j=2^{-2^j},\qquad N_j=2^{\lfloor j\nu\rfloor}\asymp2^{j\nu}.\] Use \(N_j\) separated cubes of side \(\delta_j\), with \[b_j=N_{j+1}/N_j\le2^{\lceil\nu\rceil}\] children per parent. There are enough dyadic locations for all \(j\ge j_0\); choose the initial cubes and expanded cutoff supports inside a fixed compact subset of \((0,1)^n\). Give each level-\(j\) cylinder natural mass \(1/N_j\).\par
For \(\delta_{j+1}\le r\le\delta_j\), a radius-\(r\) ball intersects only a dimensional bounded number of level-\(j\) dyadic cubes. Since adjacent gauge values have bounded ratio, \[\mu(B(x,r))\le C/N_j\le C2^{-j\nu}\le C'\varphi(r).\] The mass-distribution argument gives positive \(\mathcal H^\varphi(M)\), and level covers give finite measure. Covering a fixed cylinder by its descendants gives \[\mathcal H^\varphi(Q\cap M)\le C/N_j\le C2^{-j\nu}.\] The later proof needs these comparisons and bounded branching, not the exact integer-part count in (6.12).\par
The norm cancellation remains \[-\frac{\nu}{\sigma}+\frac{\nu}{n}
       +\frac{q-n+1}{n}=0.\] Bounded overlap also supplies the \(\ell^1\) condition required in Lemma 6.4, although p. 34 states only \(\ell^\infty\): \[\sum_Q|a_Q(x,y)|
 \le C\sum_jj^{-\delta}2^{-j\nu/\sigma}<\infty.\] Retaining the fixed comparison constant in the level-\(j+2\) separation estimate yields a kernel bound \(Cj^{\delta\sigma-\mu}2^{j\nu}\). The cylinder mass cancels its exponential factor, leaving \[C\sum_jj^{\delta\sigma-\mu}<\infty
 \quad\text{for}\quad
 1<\delta<\frac{\mu-1}{\sigma}.\] This interval is nonempty by \(\mu>1+\sigma\). The replacement construction supplies the same sharpness conclusion within the checked scope.\par
Dependency and scope. Theorem 6.2’s finite-energy estimate on p. 34 explicitly uses the cylinder-mass assertion (6.14). The required \(2^{-j\nu}\) factor cancels the kernel’s \(2^{j\nu}\), leaving the summable polynomial series. The incompatible p. 32 data do not supply a set and measure with these properties. Thus the printed proof lacks the admissible witness needed for its central sharpness conclusion, warranting Major classification. The replacement construction verifies the required estimates without validating the original one or refuting the corrected theorem.\par
\section*{Further substantive checks and limits}
For Theorem 3.1, fix \(\lambda>\|\nabla u\|_{L^A}\). The function \(A(|\nabla u|/\lambda)\) is integrable, and absolute continuity of its integral on small-volume cube unions supplies the required uniform threshold. This does not require a \(\Delta_2\) condition or order continuity of the Orlicz norm. In the content proof, choose the threshold below the cost of a fixed-size cube; then every cube in a sufficiently inexpensive cover has bounded diameter, supplying the scale restriction used in (5.14).\par
For Theorem 3.4, the decisive estimate is \[\left|\bigcup_iQ_i\right|
 \le C(\mathcal H^\varphi(E)+\varepsilon)
       \frac{\delta^n}{\varphi(\delta)}
 \longrightarrow0.\] The proof fixes its auxiliary parameter first, shrinks \(\delta\), and only afterwards sends that parameter to zero. This is consistent with the quantifiers and correctly excludes the Lebesgue endpoint from the finite-measure-to-zero assertion. The bars over \(E\) and \(E_j\) on pp. 25–26 were checked on the actual images: the argument concerns compact closures. No compactness error based on text extraction is reported.\par
The \(\Xi_{\psi,B}\) estimate in Theorem 3.5 follows by choosing \(r=\kappa/(c_n\mathcal H_\infty^\varphi(E))\) after the cover sum. The infinite-content case is treated by exhaustion; the relevant continuity-from-below theorem remains an explicit external dependency.\par
For the sharpness construction, the logarithmic cutoff has \(|\nabla\eta_j|\asymp(2^j|x|)^{-1}\) on its annulus. With \(\lambda_j=2^{j(q-n+1)/n}\), its modular integral is bounded by \[\frac{C}{2^{jn}\lambda_j^n}
 \int_{\delta_j}^{\delta_{j-1}}
 \frac{\log^q(1/r)}r\,dr\le C.\] This also works when \(n=1,q>0\). Conditioning on a random variable with maximal coefficient reduces the used positive-logarithmic kernel estimate to a radial ball integral; \(\sigma<n\) makes it finite.\par
A bounded reconstruction of the omitted proof of Theorem 6.1 uses separated cubes of side \(2^{-Lj}\), number comparable to \(2^{Lj\alpha}\), and coefficients \[j^{-\delta}2^{-Lj(\alpha+p-n)/p}.\] Their gradient norms sum for \(\delta>1\). With \(\sigma=\alpha p/(\alpha+p-n)\), the same random-energy argument leaves \(\sum_jj^{\delta\sigma-\gamma}\), convergent when \(1<\delta<(\gamma-1)/\sigma\). This verifies the exponent mechanism stated in the paper. It does not resolve the smaller logarithmic exponents or the general-gauge optimality questions that the authors explicitly leave open.\par
\section*{Disposition}
The assessment is Major because F4 invalidates a derivation used in the main distortion estimates and F7 invalidates the construction supplying the finite-energy witness for Theorem 6.2. The explicit replacements restore these steps; no counterexample to the principal distortion or sharpness conclusions is established. F1–F3, F5 and F6 retain their stated local scope and Minor classification. The external-dependency limits and the distinction between the examined manuscript and the unexamined journal version remain. Scientific decision and adjudication are pending; author adoption is unknown.\par
9 A. Cianchi, M. V. Korobkov and J. Kristensen, Distortion of Hausdorff measures under Orlicz–Sobolev maps, arXiv:2208.08152v2, 26 May 2023, 37 pages. https://arxiv.org/abs/2208.08152v2. EMS Press, article record, DOI 10.4171/JEMS/1620, https://ems.press/journals/jems/articles/14298665. Used for identity and version separation.\par
\end{document}
